% AUTORIDAD DE FUENTE:
% expose_VIII_body.tex, líneas 1745--1799; índices de escaneo 286--289;
% folios de fuente 275--278. Las cuatro correcciones materiales de la
% capa francesa se conservan reversiblemente mediante \sisi.

\medskip
\noindent 4.12. Sea $P$ un esquema en grupos liso y conexo, sea
$\mathbb L$ un haz invertible sobre $P$, y pongamos
\begin{equation}\tag{4.12.1}
 \delta(\mathbb L)=
 \pi^*\mathbb L\otimes
 \pr_1^*(\mathbb L)^{-1}\otimes
 \pr_2^*(\mathbb L)^{-1},
\end{equation}
con la notación de 4.2. Como ya se observó en VII 1.3.4, este es un haz
invertible birrigidificado sobre $P\times P$. Por 4.3, si
$\delta(\mathbb L)$ procede de una biextensión, esta queda determinada
salvo isomorfismo único. En particular, por 4.7, está definida cuando
$P$ es una extensión de un esquema abeliano $A$ por una extensión
iterada de toros y copias de $\underline G_a$. Siempre que está
definida, esta biextensión es manifiestamente simétrica.

%%%%%%%%%%  scan idx 287  |  folio 276  %%%%%%%%%%

En el caso favorable precedente, y suponiendo que el toro maximal de
$P$ es escindido, la construcción de $\delta(\mathbb L)$ puede
reducirse al caso en que $P$ es un esquema abeliano. En efecto, la
construcción es funtorial en $P$, mientras que el homomorfismo
\begin{equation}\tag{4.12.2}
 \Pic(A)\longrightarrow\Pic(P)
\end{equation}
inducido por el homomorfismo cociente canónico
\[
 f:\sisi{A\longrightarrow P}{P\longrightarrow A}
\]
es sobreyectivo por 4.5 (véase el comienzo de la demostración de 4.6).
Por tanto, podemos encontrar un isomorfismo
\[
 \mathbb L\sisi{^*}{}\simeq f^*\mathbb M,
\]
donde $\mathbb M$ es un haz invertible sobre
$\sisi{\mathbb M}{A}$. Entonces $\delta(\mathbb L)$ es la imagen
inversa de la biextensión $\delta(\mathbb M)$ de $(A,A)$ por
$\underline G_{m,k}$.

\medskip
\noindent\underline{Proposición 4.13.} \underline{Bajo las hipótesis
de 4.12, supongamos que $\mathbb L$ es un haz invertible amplio sobre
$P$. Si la biextensión $\delta(\mathbb L)$ de $(P,P)$ por
$\underline G_{m,k}$ está definida, entonces es no degenerada en el
sentido de 4.11. Si $P$ es una extensión de un esquema abeliano $A$ por
un grupo afín liso y conexo $L$ que satisface la condición de 4.7,
entonces la biextensión de $(A,A)$ de la cual procede
$\delta(\mathbb L)$ por 4.8 es de la forma $\delta(\mathbb M)$, donde
$\mathbb M$ es un haz invertible amplio sobre $A$.}

Nos reducimos al caso en que $k$ es algebraicamente cerrado. Por la
definición, basta entonces, con la notación de 4.12, demostrar que la
biextensión $\delta(\mathbb M)$ de $(A,A)$ es no degenerada. M. Raynaud
demuestra en su tesis [2, XI 1.11] que, si
$\mathbb L=f^*\mathbb M$ es amplio, entonces $\mathbb M$ también lo es.
La teoría de A. Weil da entonces la conclusión conocida de que
$\delta(\mathbb M)$ es no degenerada [1, 2], es decir, que el
homomorfismo de $A$ a su esquema abeliano dual es una isogenia.

%%%%%%%%%%  scan idx 288  |  folio 277  %%%%%%%%%%

\medskip
\noindent\underline{Observación 4.14.} Supongamos que estamos en el
caso favorable en que $P$ es una extensión de un esquema abeliano $A$
por un grupo $L$ que es una extensión iterada de toros y copias de
$\underline G_a$, y consideremos el compuesto canónico
\begin{equation}\tag{4.14.1}
 \operatorname{NS}(A)
 \xrightarrow{\ \delta\ }
 \Biext^1(A,A;\underline G_m)
 \xrightarrow{\ \sim\ }
 \Biext^1(P,P;\underline G_m).
\end{equation}
Por 4.8, la segunda flecha es biyectiva. Por otra parte, la teoría de
las variedades abelianas muestra que la primera flecha identifica el
grupo de Néron--Severi $\operatorname{NS}(A)$ de $A$ con el subgrupo de
\[
 \Biext^1(A,A;\underline G_m)=\operatorname{Corr}(A,A)
\]
formado por las clases de biextensiones simétricas o,
equivalentemente, de correspondencias divisoriales simétricas. Así,
(4.14.1) identifica $\operatorname{NS}(A)$ con el grupo de las
biextensiones simétricas de $(P,P)$ por $\underline G_m$. Por 4.7,
este grupo también puede describirse como el grupo de clases de
correspondencias divisoriales simétricas sobre $(P,P)$.

Es natural llamar \underline{amplia} a una correspondencia divisorial
sobre $(P,P)$, y llamar asimismo \underline{amplia} a la biextensión
correspondiente de $(P,P)$ por $\underline G_m$, cuando procede de un
elemento amplio de $\operatorname{NS}(A)$. Esta condición es
manifiestamente invariante por extensión del cuerpo base; por tanto, la
noción sigue teniendo sentido sin la hipótesis restrictiva precedente
sobre $P$, suponiendo solo que $P$ es liso y conexo.

%%%%%%%%%%  scan idx 289  |  folio 278  %%%%%%%%%%

Supongamos además que $L$ es ``resoluble escindido sobre $k$'', en el
sentido de que admite una dévissage por copias de $\underline G_a$ y
$\underline G_m$. La observación precedente y la sobreyectividad de
(4.12.2) muestran que la imagen de
\begin{equation}\tag{4.14.2}
 \Pic(P)\xrightarrow{\ \delta\ }
 \Biext^1(P,P;\underline G_m)
\end{equation}
está formada por las biextensiones simétricas, y que su núcleo es el
subgrupo $\Pic^0(P)$ de $\Pic(P)$ que es imagen de $\Pic^0(A)$. En
consecuencia, si ponemos
\[
 \operatorname{NS}(P)=\Pic(P)/\Pic^0(P)
 \simeq\operatorname{NS}(A),
\]
también podemos interpretar (4.14.1) como un isomorfismo canónico
\begin{equation}\tag{4.14.3}
 \operatorname{NS}(P)
 \xrightarrow{\ \sim\ }
 \Biext^1(P,P;\underline G_m)^{\mathrm{sym}}.
\end{equation}

Finalmente, el teorema de Raynaud citado en 4.13 afirma que un elemento
$\xi$ de $\Pic(P)$ es amplio si y solo si su imagen $\delta(\xi)$ por
$(\sisi{4.14.1}{4.14.2})$ es la clase de una biextensión amplia.
